\documentclass[11pt]{article}
\usepackage[margin=1in]{geometry}
\usepackage{amsmath,amssymb,amsthm,hyperref}
\newtheorem{theorem}{Theorem}
\newtheorem{lemma}{Lemma}
\newtheorem{proposition}{Proposition}
\newcommand{\E}{\mathbb E}
\newcommand{\Prb}{\mathbb P}
\newcommand{\Var}{\operatorname{Var}}
\newcommand{\Poi}{\operatorname{Poi}}
\newcommand{\supp}{\operatorname{supp}}
\title{Reducing the explicit sampling constant in Optimal Support Size Testing}
\author{Research derivation for the supplied manuscript; not externally peer reviewed}
\date{21 October 2026}
\begin{document}
\maketitle

\paragraph{Scope.}
We analyze the exact real-parameter, Poissonized algorithm in Section 3 of the
supplied manuscript by Qisheng Wang, \emph{Optimal Support Size Testing},
local file \texttt{2610.18723v1.pdf}. We keep $L_0=1024$, the coefficient $64$
in $b$, every scale, every polynomial degree, and every decision correction
unchanged. Only $C_0$ is reduced. The bounds replacing Lemma 4.1 below are new
derivations; the original intermediate bound $\|f\|_\infty\le4n^{1/16}$ is not
asserted with the new constants. The error conclusions of Lemma 4.2 and the
main theorems are preserved. This is not a deduction from experimental curves.

\begin{theorem}[Sufficient constant]
Set
\[
 \eta=\frac{9}{4\cdot10^{219}}=2.25\cdot10^{-219},\qquad
 C_* =1024+\eta.
\]
For every $C\in[C_*,2^{26}]$, the original Section 3 algorithm with $C_0=C$
and $L_0=1024$ satisfies, for every integer $n\ge1$ and $0<\epsilon<1$,
\begin{align}
 S=|\supp(p)|\le n&\ \Longrightarrow\
 \Prb(\widetilde S>S)<0.01,\label{eq:upper}\\
 1\le r\le n,\quad\Delta_r(p)\ge\epsilon&\ \Longrightarrow\
 \Prb(\widetilde S\le r)<0.01.\label{eq:lower}
\end{align}
Consequently the tester has error below $1\%$ on each promise class, and
$\widehat S=\min\{\max\{\lceil\widetilde S\rceil,1\},n\}$ satisfies
\[
 \Prb\{\min(\operatorname{eff}_\epsilon(p),n)\le\widehat S\le S\}\ge0.98.
\]
In particular the simple integer choice $C_0=1025$ is sufficient. Its
sampling coefficient is smaller than the original by a factor
$2^{26}/1025>65{,}472$.
\end{theorem}

\paragraph{Integer rounding of the mean.}
For either recommended choice $C=C_*$ or $C=1025$, replacing $m$ by
$\lceil Cn/(\epsilon L)\rceil$ and recomputing $b$ preserves the result.
In the count branch this only increases the sample mean. In the polynomial
branch it amounts to using an effective coefficient
$C'\in[C,C+\epsilon L/n)$, and $L/n\le1024e^{-1023}<1$ for $L\ge1024$.
Thus $C'\in[C_*,2^{26}]$, where all the pointwise bounds below apply.

\begin{theorem}[Obstruction with the unchanged count branch]
For every $0<C\le1024$, there exist a finite distribution, an integer $n$,
and $\epsilon=1/4$ for which $L<1024$, $\Delta_n(p)=\epsilon$, and the
unchanged Section 3 tester with mean $m=Cn/(\epsilon L)$ has error above
$1\%$. Thus no such $C$ preserves the original $1\%$ guarantee.
Together with the preceding theorem, the admissible constant is localized
to within $2.25\cdot10^{-219}$ of the obstruction at $1024$.
This does not identify the exact optimum in that tiny interval, nor is the
obstruction a lower bound for all support testers or for a changed $L_0$.
\end{theorem}

\paragraph{Integer optimum under the stated guarantee.}
The two theorems show that $1025$ is the smallest positive integer $C_0$
preserving the original $1\%$ guarantee with the rest of Section 3 fixed.
This integer-optimality statement is specifically about that stronger
guarantee, not about all algorithms achieving success probability $2/3$.

\section{Algorithm and notation}
All logarithms except $\log_2$ are natural. Put
\[
 L=1+\ln n,\quad m=\frac{Cn}{\epsilon L},\quad
 M\sim\Poi(m),\quad N_i\sim\Poi(mp_i)\ \text{independently},\quad
 D=\sum_i\mathbf1\{N_i>0\}.
\]
If $L<1024$, set $\widetilde S=D$. Otherwise define
\[
 a_0=\frac{\epsilon}{8n},\quad b=\frac{64L}{m},\quad
 J=\left\lfloor\log_2\frac b{2a_0}\right\rfloor,\quad
 a_j=2^ja_0,\quad \delta_j=2^{-j-2},\quad
 k_j=\left\lceil\sqrt{b/a_j}\ln(2/\delta_j)\right\rceil.
\]
For each scale, write
\[
 P_j(x)=-\rho_jT_{k_j}\!\left(\frac{b+a_j-2x}{b-a_j}\right)
       =\sum_{\ell=0}^{k_j}c_{j,\ell}x^\ell,\qquad
 \rho_j=T_{k_j}\!\left(\frac{b+a_j}{b-a_j}\right)^{-1}.
\]
Use $f_j(0)=0$, $f_j(\ell)=1+c_{j,\ell}\ell!/m^\ell$ for
$1\le\ell\le k_j$, and $f_j(\ell)=1$ above $k_j$. Finally
\[
 Z_j=\sum_i f_j(N_i),\qquad
 \widetilde S=\max\left\{D,\max_{0\le j\le J}\frac{Z_j-\delta_jn}{1+\delta_j}\right\}.
\]
The output is Yes iff $\widetilde S\le n$. These are exactly the original
formulas with the revised $C$; $m$ is a Poisson mean, not a realized sample
count. A numerical implementation must not round $1024+\eta$ to $1024$.

\section{A sharp enough count bound}
\begin{lemma}
If $\Delta_r(p)\ge\epsilon$, then for the Poissonized observed count,
\[
 \Prb(D\le r)\le\Prb\{\Poi(m\epsilon)\le r\}.
\]
\end{lemma}
\begin{proof}
Embed sampling in a rate-one marked Poisson process, with marks drawn from
$p$. After any $k\le r$ distinct labels have been seen, their total mass is
at most that of the $r$ largest categories. Thus the rate of the next new
label is at least $\epsilon$. Conditional on the labels already seen, the
waiting time for the next new label is exponential with that rate.
The independent exponential-race construction couples each of the first
$r+1$ waiting times to be at most an independent $\mathrm{Exp}(\epsilon)$
waiting time. Hence the time to discover $r+1$ labels is stochastically
at most $\mathrm{Gamma}(r+1,\epsilon)$. Its survival probability at time $m$
is $\Prb\{\Poi(m\epsilon)\le r\}$.
\end{proof}

For $x>1$, exponential Markov applied with $s=\ln x$ gives
\begin{equation}
 \Prb\{\Poi(nx)\le n\}
 \le \exp\{-n g(x)\},\qquad g(x)=x-1-\ln x.
 \label{eq:chernoff}
\end{equation}
Here $g$ is increasing on $(1,\infty)$, and, for $h\ge0$,
\begin{equation}
 g(1+h)=\int_0^h\frac{u}{1+u}\,du
 \ge\frac{h^2}{2(1+h)}.\label{eq:g}
\end{equation}

For the small-$L$ branch, $x=C/L>1$. Set $t=1022.999$ and
$h=\eta/1024$. We check all integers $n$ using three ranges:
\begin{enumerate}
\item If $n\le e^{40}$, then $L\le41$ and
\[
 ng(C/L)\ge g(1024/41)>20.75>\ln100.
\]
\item If $e^{40}<n\le e^t$, then $L\le t+1=1023.999$ and
$C/L\ge1+h_2$, where $h_2=1024/1023.999-1=1/1023999$. Therefore
\[
 ng(C/L)>\frac{e^{40}h_2^2}{2(1+h_2)}>112240>\ln100.
\]
\item If $e^t<n<e^{1023}$, then $L<1024$ and $C/L>1+h$, giving
\[
 ng(C/L)>\frac{e^t h^2}{2(1+h)}>4.6297>\ln100.
\]
\end{enumerate}
Thus, if $r\le n$ and $\Delta_r(p)\ge\epsilon$,
\[
 \Prb(\widetilde S\le r)\le\Prb\{\Poi(m\epsilon)\le n\}<0.01.
\]
Also $D\le S$ surely. This proves both desired conclusions for $L<1024$.
The decimal inequalities above are strict bounds on explicit elementary
expressions, independently checked by directed interval arithmetic in the
accompanying script; they are not statistical estimates.

\section{Parameter bounds in the polynomial branch}
Suppose $L\ge1024$ and $1024\le C\le2^{26}$. Then
\[
 \frac b{2a_0}=\frac{256L^2}{C}\ge4,
 \qquad 0<a_j\le b/2,\qquad
 b=\frac{64\epsilon L^2}{Cn}\le\frac{L^2}{16e^{L-1}}<1.
\]
The last function is decreasing for $L>2$, and its value at $1024$ is
less than $4\cdot10^{-440}$. Let $t_j=a_j/b$. Then
\[
 t_j=\frac{C2^j}{512L^2},\qquad
 k_j\le \sqrt{\frac{512}{C}}L\frac{(j+3)\ln2}{2^{j/2}}+1.
\]
The sequence $(j+3)2^{-j/2}$ is decreasing for integers $j\ge0$, since
$(j+4)/(j+3)\le4/3<\sqrt2$. Consequently
\begin{equation}
 \frac{k_j}{L}\le\frac{3\ln2}{\sqrt2}+\frac1{1024}
 <1.472<\frac32.\label{eq:degree}
\end{equation}
In particular
\[
 m\sqrt{b(b-a_j)}\ge\frac{64L}{\sqrt2}>k_j,
 \qquad m(b-a_j)\ge32L.
\]
Furthermore,
\[
 \rho_j=\frac1{\cosh(2k_j\operatorname{artanh}\sqrt{a_j/b})}
 \le2e^{-2k_j\sqrt{a_j/b}}
 \le\frac{\delta_j^2}{2}\le\delta_j.
\]

\section{Expectation bound, with its analytic prerequisites checked}
For a general scale $a,b,k,m$, direct substitution into the Poisson series
gives
\[
 Q(x)=\E[f(\Poi(mx))]=1+e^{-mx}P(x).
\]
We recall a short proof of the bound used in the supplied Lemma 2.3.
On $[0,a]$, $P$ is nonpositive, increasing and concave: $T_k$ is increasing
and convex on $[1,\infty)$. Since $P(0)=-1$ and $P(a)=-\rho$, its chord gives
$Q(x)\ge(1-\rho)x/a$, while $Q(x)\le1$. On $[a,b]$, $|P(x)|\le\rho$, so
$1-\rho\le Q(x)\le1+\rho$. For $x\ge b$, put $u=2(x-b)/(b-a)$.
Using $T_k(1+u)=\cosh(k\operatorname{arcosh}(1+u))\le e^{k\sqrt{2u}}$
and $x^2\ge4b(x-b)$, the condition $m\sqrt{b(b-a)}\ge k$ implies
$mx\ge k\sqrt{2u}$. Hence $|Q(x)-1|\le\rho$ also in this range.
Combining the three ranges yields
\begin{equation}
 (1-\delta_j)\min\{x/a_j,1\}\le Q_j(x)\le1+\delta_j,
 \qquad x\ge0.\label{eq:expectation}
\end{equation}
In particular
\[
 (1-\delta_j)\sum_i\min\{p_i/a_j,1\}\le\E Z_j\le(1+\delta_j)S.
\]

\section{An improved bound on polynomial weights}
\begin{lemma}
For $0<a<b$, $k\ge1$, and the original Chebyshev weights,
\[
 \|f\|_\infty\le4\exp\!\left(\frac{2k^2}{m(b-a)}\right).
\]
\end{lemma}
\begin{proof}
Let $z_0=(b+a)/(b-a)>1$. All roots of $T_k$ are in $(-1,1)$ and its
leading coefficient is positive, so every coefficient of $T_k(z_0+v)$
as a polynomial in $v$ is positive. Therefore, for $u\ge0$,
\[
 \sum_{\ell=0}^k|c_\ell|u^\ell
 =\rho T_k\!\left(z_0+\frac{2u}{b-a}\right),\qquad |c_0|=1.
\]
By $\int_0^\infty e^{-s}s^\ell ds=\ell!$,
\begin{equation}
 A:=\sum_{\ell=0}^k\frac{|c_\ell|\ell!}{m^\ell}
 =\int_0^\infty e^{-s}
   \frac{T_k(z_0+2s/[m(b-a)])}{T_k(z_0)}\,ds,
 \qquad \|f\|_\infty\le A.\label{eq:integral}
\end{equation}
To bound the ratio, write $v=a/(b-a)$ and $u=s/[m(b-a)]$.
The identity $\operatorname{arcosh}(1+2v)=2\operatorname{arsinh}\sqrt v$
and the fact that $\operatorname{arsinh}$ is 1-Lipschitz on $[0,\infty)$
imply
\[
 \operatorname{arcosh}(1+2v+2u)-\operatorname{arcosh}(1+2v)
 \le2(\sqrt{v+u}-\sqrt v)\le2\sqrt u.
\]
Since $T_k(z)=\cosh(k\operatorname{arcosh}z)$,
\[
 \frac{T_k(z_0+2u)}{T_k(z_0)}\le2e^{2k\sqrt u}.
\]
(The factor 2 can be improved to 1, but is not needed.) Now
\[
 2k\sqrt{\frac{s}{m(b-a)}}\le\frac s2+\frac{2k^2}{m(b-a)}.
\]
Substituting in \eqref{eq:integral} and integrating $e^{-s/2}$ gives
$A\le4\exp(2k^2/[m(b-a)])$.
\end{proof}

By \eqref{eq:degree}, at every scale
\begin{equation}
 \|f_j\|_\infty\le4e^{9L/64}.\label{eq:weights}
\end{equation}
This replaces the loose coefficient bound $1+15^{k_j}$ in the original
proof. It is the main reason the huge original $C_0$ is unnecessary.

\section{Simultaneous concentration}
Because $f_j(0)=0$ and the category counts are independent,
\[
 \Var(Z_j)\le\|f_j\|_\infty^2\E D.
\]
If $\E D\le2n$, then \eqref{eq:weights} and Chebyshev give
\[
 \Prb\{|Z_j-\E Z_j|>\delta_jn\}
 \le\frac{32e^{9L/32}}{n\delta_j^2}.
\]
Also
\[
 \delta_J^{-1}=4\cdot2^J\le\frac{1024L^2}{C}\le L^2,
 \qquad \sum_{j=0}^J\delta_j^{-2}\le\frac43\delta_J^{-2}\le\frac43L^4.
\]
Thus the event
\[
 \mathcal E=\{|Z_j-\E Z_j|\le\delta_jn\text{ for every }j\}
\]
satisfies
\begin{equation}
 \Prb(\mathcal E^c)
 \le\frac{128e}{3}L^4e^{-23L/32}
 \le\frac{128e}{3}(1024)^4e^{-736}
 <3\cdot10^{-306}<0.01.\label{eq:concentration}
\end{equation}
The function decreases for $L\ge1024$, since its logarithmic derivative
is $4/L-23/32<0$. For countable alphabets the variance calculation follows
by finite truncation: $\E D\le m<\infty$ and bounded weights ensure the
required $L^2$ convergence.

\section{Completion of the correctness proof}
If $S\le n$, then $\E D\le S\le n$ and, on $\mathcal E$,
\[
 \frac{Z_j-\delta_jn}{1+\delta_j}
 \le\frac{\E Z_j}{1+\delta_j}\le S.
\]
Together with $D\le S$, this proves \eqref{eq:upper}.

Now fix $1\le r\le n$ with $\Delta_r(p)\ge\epsilon$.
If $\mu=\E D>2n$, independence of the occupancy indicators gives
$\Var(D)\le\mu$. Therefore
\[
 \Prb(\widetilde S\le r)\le\Prb(D\le n)
 \le\frac{\mu}{(\mu-n)^2}\le\frac2n<0.01.
\]
Otherwise $\E D\le2n$ and \eqref{eq:concentration} applies. Order the
probabilities as $p_{(1)}\ge p_{(2)}\ge\cdots$ and consider three cases.

\paragraph{Case 1: $p_{(r)}\ge b/2$.}
If the top $r$ labels and at least one tail label are observed, $D>r$.
A union bound and Poisson zero probabilities give
\[
 \Prb(D\le r)\le r e^{-mb/2}+e^{-m\Delta_r(p)}
 \le ne^{-32L}+e^{-Cn/L}<0.01.
\]
Indeed $ne^{-32L}=e^{-31L-1}$, while $Cn/L\ge1024e^{L-1}/L$,
whose right side is increasing for $L\ge1024$ and at $1024$ equals $e^{1023}$.

\paragraph{Case 2: $p_{(r)}<a_0$.}
All tail weights are below $a_0$, so
\[
 \E Z_0\ge\frac{1-\delta_0}{a_0}\Delta_r(p)\ge6n.
\]
On $\mathcal E$,
\[
 \widetilde S\ge\frac{\E Z_0-2\delta_0 n}{1+\delta_0}
 \ge\frac{6n-n/2}{5/4}=\frac{22}{5}n>r.
\]

\paragraph{Case 3: $a_0\le p_{(r)}<b/2$.}
There is $0\le j\le J$ such that $a_j\le p_{(r)}<2a_j$.
For every tail weight, $\min\{p_i/a_j,1\}\ge p_i/p_{(r)}$, and the top
$r$ summands equal one. Consequently
\[
 \E Z_j\ge(1-\delta_j)\left(r+\frac{\epsilon}{2a_j}\right).
\]
As $a_j\delta_j=\epsilon/(32n)$, on $\mathcal E$,
\begin{align*}
 \widetilde S
 &\ge\frac{\E Z_j-2\delta_jn}{1+\delta_j}\\
 &\ge r+\frac{\delta_j\{(14-16\delta_j)n-2r\}}{1+\delta_j}>r,
\end{align*}
because $\delta_j\le1/4$ and $r\le n$.
These cases prove \eqref{eq:lower}. Their probabilities are not added:
the case selection depends on the fixed distribution, and only one applies.

For the lower-support output, the upper inequality is automatic if $S\ge n$;
otherwise \eqref{eq:upper} and integer-valued $S$ give it with probability
at least $0.99$. Put $r=\min\{\operatorname{eff}_\epsilon(p),n\}-1$.
If $r=0$, the lower inequality is automatic. Otherwise
$\Delta_r(p)>\epsilon$ by the definition of effective support, and
\eqref{eq:lower} implies $\widehat S\ge r+1$ with probability at least $0.99$.
The union bound proves joint success at least $0.98$.

The sample order remains $\Theta(n/[\epsilon\log(en)])$.
For a deterministic sample cap, abort when $M>4m$ and allow an arbitrary
output on that event. The Poisson tail adds at most $e^{-m}\le e^{-1024}$
to the failure probability and uses at most $\lceil4m\rceil$ samples.
Thus the original $2/3$-success main theorems in the ordinary sampling model
also follow. The optimized $C$ itself is a Poisson-mean coefficient;
de-Poissonization introduces its own constant factor.

\section{Why $C\le1024$ cannot preserve the 1\% guarantee}
Let $n=\lceil e^{1023}\rceil-1$, so $1023<L=1+\ln n<1024$ and the
algorithm takes the count branch. At $C=1024$, write
\[
 \lambda=\frac{1024n}{L}=n+d,\qquad
 0<d=\frac{n\ln(e^{1023}/n)}L\le\frac1L<\frac1{1023}.
\]
Here we used $e^{1023}-n\le1$ and $\ln(1+x)\le x$.

Let $Y\sim\Poi(n)$ and $Z=n-Y$. The Poisson central moments are
$\E Z=0$, $\E Z^2=n$, and $\E Z^4=3n^2+n\le4n^2$.
H\"older's inequality gives
\[
 \E|Z|\ge\frac{(\E Z^2)^{3/2}}{(\E Z^4)^{1/2}}\ge\frac{\sqrt n}{2}.
\]
Since $\E Z_+=\E|Z|/2$, Cauchy--Schwarz yields
\[
 \Prb(Y<n)=\Prb(Z>0)
 \ge\frac{(\E Z_+)^2}{\E Z_+^2}\ge\frac1{16}.
\]
Couple an independent $V\sim\Poi(d)$ to obtain
\begin{equation}
 \Prb\{\Poi(n+d)\le n-1\}
 \ge\Prb(V=0)\Prb(Y<n)
 \ge\frac{e^{-d}}{16}>\frac{e^{-1/1023}}{16}>0.0624.
 \label{eq:obstruction}
\end{equation}
This bound remains true when $C$ is smaller, by monotonicity in the mean.

To realize this in a finite promised distribution, fix $\epsilon=1/4$.
Choose a large integer $K$ and put
\[
 q=\frac{\epsilon}{K-n+1},\qquad
 p_0=1-Kq,\qquad p_i=q\quad(1\le i\le K).
\]
For sufficiently large $K$ the head $p_0$ is positive and is the largest
atom, and
\[
 \Delta_n(p)=(K-n+1)q=\epsilon.
\]
The total number of tail observations is $Y_K\sim\Poi(mKq)$. If
$Y_K\le n-1$, then $D\le Y_K+1\le n$, so the tester incorrectly returns
Yes. As $K\to\infty$, $mKq\to m\epsilon=Cn/L$, and the finite Poisson
tail sum is continuous in its mean. By \eqref{eq:obstruction}, for a
sufficiently large finite $K$ the error exceeds $0.01$ (in fact $0.06$).
This proves the obstruction.
If the mean is rounded upward, the tail mean increases by less than
$\epsilon=1/4$. The same argument gives the still sufficient lower limit
$e^{-(1/1023+1/4)}/16>0.048$, so the obstruction is not a rounding artifact.

\paragraph{What is and is not optimized.}
The result preserves the stronger $1\%$ conclusion of the original proof,
not just the $2/3$ statement. It brackets the constant for that guarantee
with the original branch threshold fixed. The lower bound above alone does
not prove failure of a $2/3$-success guarantee at $C=1024$; its certified
failure probability is only above $6\%$. Nor does it rule out much smaller
constants after changing the branch threshold or the algorithm. No claim
of the exact optimal increment above 1024, or of a new asymptotic sample
complexity, is made. The easy-to-use integer $1025$ avoids the extreme
precision required to encode the sharper mathematical value.

\paragraph{Verification and provenance.}
The proof uses the supplied paper's definitions, and re-proves the analytic
expectation property needed from its Lemma 2.3. The standard Chebyshev
identity can also be checked at NIST DLMF,
\url{https://dlmf.nist.gov/18.5.E1}. The companion
\texttt{verify\_c0\_reduction.py} evaluates the scalar endpoints using
4096-bit Arb interval arithmetic and exact rational constants. It checks
numerical inequalities, not the entire mathematical proof. Existing
experimental code and results were not changed to use the new constants.
An independent mathematical review remains appropriate before publication.
\end{document}
